\section{The finite dyadic procedure and its budgets}
\label{sec:procedure}

We seek a bound for modulated sums whose coefficients are formed from one
input bit and emitted with a later output bit. The bound must be uniform
in the tree depth and in the number of frequencies. We state it first;
the construction in this section supplies the pieces and budgets used in
its proof in \Cref{sec:assembly}.

\begin{proposition}[Modulated delayed Haar estimate]\label{prop:dyadic}
There are absolute constants $C$ and $d_0$ with the following property.
Let $d\ge d_0$, put $h=\log\log d$, and let $I\subset\R$ be a half-open
interval with $1/2\le |I|\le1$.  Let $\rho:I\to[0,\infty)$ be measurable,
with $\rho\le e^d$ almost everywhere and
\[
 m_{\rm top}=\frac1{|I|}\int_I\rho\le2.
\]
Give $I$ its normalized uniform probability $\mathbf U$.  Write
$\mathcal F_j$ for its depth-$j$ dyadic sigma field, $p_j(x)$ for the
depth-$j$ interval containing $x$, and $r_{j+1}$ for the sign that is $+1$
on the left child and $-1$ on the right child of each such interval.
Let $J\ge1$ and $Q\ge0$ be integers.  For $0\le j<J$, let
$\epsilon_j:I\to\C$ be $\mathcal F_{j+1}$-measurable with
$|\epsilon_j|\le1$, and use the same $\epsilon_j$ for all modulations.
The bits through depth $J+1$ are retained, so the last output sign is
$r_{J+1}$.

Set
\[
 \tau(x)=\min\bigl(\{0\le j<J:\ \E_{p_j(x)}\rho>e^{4h}\}\cup\{J\}\bigr),
 \qquad \E_p v=\frac1{|p|}\int_p v(y)\,dy,
\]
and, with $\chi_u(y)=e^{2\pi iuy}$, define
\[
 \mathcal H_{I,J}^{u}\rho(x)
 =\sum_{0\le j<\tau(x)}
       \E_{p_j(x)}[\rho(Y)\chi_u(Y)r_{j+1}(Y)]\,
       \epsilon_j(x)r_{j+2}(x).
\]
Then
\begin{equation}\label{eq:dyadic-estimate}
 \int_I\max_{\substack{u\in\Z\\|u|\le Q}}
           |\mathcal H_{I,J}^{u}\rho(x)|\,\frac{dx}{|I|}
 \le Cd.
\end{equation}
The constants are independent of $I,\rho,J,Q$, and the multipliers.
Consequently the same estimate holds with $\sup_{u\in\Z}$ in place of
the finite maximum.
\end{proposition}

\subsection{Input pieces and the order of processing}

Fix the data of \Cref{prop:dyadic}. At an output point $x$ in a depth-$j$
interval $p$, the summand in that proposition is
\begin{equation}
 \mathbb E_p[\rho(Y)\chi_u(Y)r_{j+1}(Y)]\,
       \epsilon_j(x)r_{j+2}(x),\qquad u\in\mathbb Z.
 \label{eq:source-11}
\end{equation}
The input variable $Y$ and output point $x$ play different roles. The
coefficient is known before the output bit $r_{j+1}(x)$, whereas its
multiplier $\epsilon_j(x)$ need only be known before $r_{j+2}(x)$.
The bits retained through depth $J+1$ supply this delayed sign even for
the last contributing jump. The same multipliers are used for every piece
and every label coordinate.

Every piece and frequency uses the common stopping rule $\tau$ of
\Cref{prop:dyadic}, defined from the original density $\rho$. No processing
is needed at a node where this rule has already stopped the output path.
At all other nodes, the processing below is completed before the jump.

Discard the zero input.  If $0<m_{\rm top}<1/h$, put the whole input
into static handling at the top.  Otherwise define the true input law
\[
 d\mathbf P(Y)=\frac{\rho(Y)}{m_{\rm top}}\,d\mathbf U(Y).
\]
The word ``true'' will always refer to this input law, or to its
conditional laws, and not to uniform output sampling.

A processing state at a spatial interval $p$ consists of a subdensity
$a$, its mean $m=\mathbb E_p a$, and its recorded spatial and membership
history.  The subdensity is the restriction of $\rho$ by the membership
sets on that history.  A zero-mass state is discarded.  A live state
also carries a main thermal form $A$, a count $n$ of successful main
additions, and the charge
\[
 \mathcal Q=\sum_{\text{regular successes so far}}t.
\]
Initially $A=0$, $n=0$, and $\mathcal Q=0$.  These three quantities
carry along the entire branch and are never reset at a new epoch or
plateau.

Fix constants $R_0,C_0,C_*$ in the order specified below.  The
processing rules are as follows.
\begin{enumerate}[label=\arabic*.,leftmargin=*]
\item At every incoming live state, transfer its entire current
subdensity to static handling if
\[
 m<e^{-d},\qquad \mathcal Q>R_0d,
       \qquad\text{or}\qquad n\ge\lfloor C_0h\rfloor.
\]
This transfer ends its live processing.  Its density is subsequently
only restricted to spatial descendants; it is never split by another
membership test.

\item The initial live piece starts an epoch.  Every successful main
addition starts a new epoch on its success piece unless that piece
transfers immediately.  Within an epoch, a plateau starts with
normalization $\alpha=m$ at entry.  At a subsequent incoming state it
ends and a new plateau starts if $m>2\alpha$.  The new normalization
is the incoming mean.  No new main label is drawn at such a reset.

\item At a live state that has passed the transfer and reset rules,
apply the spectral test of \Cref{lemma:spectral-update} to
$(a,p,\alpha,A)$, using the prospective child moments.  Thus
$e^{-d}\le m\le2\alpha$.  The tests range over all computation
levels $b\in[b_0/8,1]$ that are powers of two; stabilization means
that at each such level there is no separated high collection of
size $\lceil e^{k_b}\rceil$.  If a test requests an update, its success
set $E$ is a measurable set of inputs, determined from this state.
Replace the current density by the two subdensities
\[
 a\mathbf1_E\quad\hbox{and}\quad a\mathbf1_{E^c}.
\]
Only the success piece receives the new form $A^+$ and the increment
of the count; a regular success also adds $t$ to $\mathcal Q$.
Its mean is $\delta m$, where $\delta$ is its conditional success
probability.  The failure piece, if it has positive mass, has mean
$(1-\delta)m$ and retains the old epoch, plateau, form, and charge.
For a cheap update $\delta=1$.

\item Process both resulting pieces at the same spatial interval
$p$, enforcing transfer and reset rules anew.  Continue until every
remaining live piece at $p$ passes all spectral tests.  Then, and only
then, take the spatial jump for each surviving live piece and each
static piece, and proceed to its children subject to the common stop.
\end{enumerate}

Static pieces also have plateaus: start with normalization equal to
their mean at transfer (or at initial static entry), and restart
whenever the incoming mean exceeds twice that normalization.  They
undergo no spectral tests or membership splits.  Their estimates will
use only the scalar construction.

Every test, witness, and membership set is chosen from the spatial
prefix and the preceding membership outcomes.  In particular, none of
these choices uses a realized noisy observation or the next spatial
bit.  The label experiments are attached to these already specified
discrete branches.  Start with a type-$0$ main observation, independent
of the input.  Draw a fresh main observation of type $A^+$ on every
success.  On a regular success also draw the old-type evidence clones
specified in \Cref{prop:thermal-evidence}.  Make these draws even when
that success immediately causes a transfer.  Moving to a smaller
interval, following a failure, or resetting $\alpha$ does not redraw
the main observation.  Given the input and its discrete branch, every
draw uses fresh independent error with its specified channel law.

At a jumping node the density belonging to a plateau is its final
\emph{post-processing survivor}.  Its prospective child densities
are the restrictions of that survivor, before any processing at the
children.  Removed mass has instead been assigned to another plateau
or to static handling.  Therefore the densities of all jumping pieces
at $p$ sum exactly to $\rho|_p$, and their normalized contributions, after
multiplication by their respective normalizations, sum exactly to
\Cref{eq:source-11}.  A true input follows just one of these pieces;
uniform output estimates must sum all of their starting weights.

\subsection{Statement of the budgets}

For a live plateau $e$, write $p_e$ for its root and $\alpha_e$ for its
initial mean, and put $\mathbf U_e=\mathbf U(\,\cdot\mid p_e)$.
At a final live jumping state $p$ belonging to $e$, use its post
density $a$ and its fixed main channel to define
\[
 M_\pm(z)=\frac{\mathbb E_{p_\pm}[a(Y)K_A(z-Y)]}{\alpha_e},\qquad
 M=\frac{M_++M_-}{2},\qquad
 D=\frac{M_+-M_-}{2},\qquad
 J_p=\int\frac{|D|^2}{M}\,dz.
\]
The piece and plateau are implicit in $J_p$; zero divided by zero is
taken to be zero.  Notice that $\int M=m_p/\alpha_e\le2$, where
$m_p=\mathbb E_p a$, whereas $\int M_\pm\le4$.  In particular
$J_p\le2$.

\begin{proposition}[Live budgets]\label{prop:live-budgets}
There are fixed absolute choices of the constants in the procedure
and in \Cref{lemma:spectral-update} such that, for all sufficiently
large $d$, the procedure is finite in every finite spatial tree and
preserves the exact splitting identity above.  When
$1/h\le m_{\rm top}\le2$, it has the following bounds, with constants
independent of $J$, $I$, the input, and any frequency cutoff.
\begin{enumerate}[label=\textup{(\roman*)},leftmargin=*]
\item Before every requested live update the count and form
hypotheses of \Cref{lemma:spectral-update} hold.  On every input
branch, including its final success if a cap is crossed,
\[
 \sum_{\rm successes}k_b+\sum_{\rm successes}t
       +\sum_{\rm successes}\bigl(F(A^+)-F(A+S_p)\bigr)\le Cd,
 \qquad n\le C h.
\]
The count flag cannot cause transfer unless the charge flag has
already been crossed at that success.

\item Let $B_{\rm mem}$ be the sum of the negative logarithms of the
conditional probabilities of all observed membership outcomes along
the true path through live termination.  Let
$X_* =\max\log^-m$ over its processing states, including the state
at transfer or at the common stop, but excluding later static
descendants.  Then
\begin{equation}
 \|B_{\rm mem}\|_{L^2(\mathbf P)}
       +\|X_*\|_{L^2(\mathbf P)}\le C h,
 \label{eq:source-12}
\end{equation}
and $\mathbf P(\mathrm{transfer})\le C h^{-10}$.

\item The live plateau weights satisfy
\[
 \sum_{e\text{ live}}\mathbf U(p_e)\alpha_e\le C h.
\]
If $\xi_p$ is the local average of mass removed from the survivor of
one plateau at $p$, divided by its $\alpha_e$, then
$\mathbb E_{\mathbf U_e}\sum_p\xi_p\le1$.  Full removal at an exit
may be included.  Starting from a final jumping post state $p_0$,
the corresponding future bound under uniform probability on $p_0$
is at most $m_{p_0}/\alpha_e\le2$.

\item The global label Fisher sum obeys
\begin{equation}
 \sum_{\text{all live jumping pieces }p}
        \mathbf U(p)\alpha_e J_p\le C d.
 \label{eq:source-13}
\end{equation}
\item Ignore the labels and use the filtration of spatial bits and
memberships.  After live termination one may keep the final piece
fixed and reveal any finite number of additional spatial bits, making
no further membership observations.  The expected sum of the spatial
bit relative entropies to a fair bit is $O(d)$.  If this spatial
extension is stopped at the common density cutoff, the bound is
$O(h)$.
\end{enumerate}
\end{proposition}

\begin{proof}
\subsubsection*{Caps, constant order, and finite termination}
Let $C_1$ be the numerical constant bounding a relative metric cost
by $C_1t$ in \Cref{lemma:spectral-update}, and let $C_2$ be its
numerical regular gap-increase constant.  Choose the numerical
$\varepsilon>0$ there small enough for the cheap gap decrease, and
take $R_0=1000$.  Consider a prefix of successes for which the update
hypotheses have been justified.  The gap potential
\[
 F(A)-\log|B_A(r)|
\]
starts at zero and is bounded below by an absolute constant.  Here
$r$ begins at $r_{\rm sep}/4$ and increases by
$r_{\rm sep}/(10C_0h)$ at each addition, so $r\le r_{\rm sep}/2$
under the count cap.  Telescoping its regular increases and cheap
decreases gives
\begin{equation}
 \sum_{\rm cheap}k_b\le
          2C_2\sum_{\rm regular}t+C,
 \qquad
 \sum_{\rm successes}t\le
       (1+2\varepsilon C_2)\sum_{\rm regular}t+C.
 \label{eq:06-gap-accounting}
\end{equation}
In the second inequality we used $t=\varepsilon k_b$ on cheap
successes.  Regular successes also satisfy $\varepsilon k_b\le t$.
Consequently, while $\mathcal Q\le R_0d$, the sum of all success
$k_b$'s is at most $C(R_0,\varepsilon,C_2)d$.  Since every
$k_b\ge d/h$, the number of those successes is at most $C'h$.
Choose $C_0>C'+2$.  A first count breach without a charge breach
would obey this same prefix bound, contradicting that choice for
large $d$.

The form cap is also inductive.  If the preceding main addition
occurred at $p_*$, the spatial and metric estimates give
\[
 F(A^+)-F(A)
       \le\log(R_p/R_*)+C+C_1t.
\]
The spatial scale ratios telescope along a branch.  Before a charge
breach, \Cref{eq:06-gap-accounting} and the count estimate therefore
give $F(A)\le\log R_*+C_*d$ for a sufficiently large fixed $C_*$.
The initial case is $A=0$ with $p_*=I$.  At the next update, the
flat-increment condition \Cref{eq:source-5} is supplied by the update
itself.  Thus the assumptions needed to construct each successive
update have been established from its predecessors.

Fix $D_0,P_0$ in the compression application using these caps, then
fix $A_0$ sufficiently large for the cardinality loss, and finally
take $d$ sufficiently large for all required asymptotic inequalities.
The constants $C_1,C_2$ and the choice of $\varepsilon$ precede the
caps; their numerical nature is what permits this order.  A regular
success that first crosses the charge cap adds at most $D_0d$ and
is followed immediately by transfer.  Hence that final overshoot
only changes the absolute constants in the claimed path budgets.
No subsequent update uses a form cap for this terminal form.  The
energy estimate in \Cref{eq:source-7} remains valid for the labels
drawn at that success, as it uses the flat-increment and count bounds.
This proves all the deterministic budget assertions.

At a fixed spatial node, the number of successes along any recurrence
path is bounded by the count cap.  Between successes, every failure
multiplies $m$ by at most $1-c(b_0/8)^2$.  Since $m\le e^d$, only
finitely many consecutive failures are possible before $m<e^{-d}$.
There is a uniform finite bound, depending on $d$ but not the node,
on the length of this local binary recursion.  Its branching is
finite; induction over the finite spatial depth proves termination of
the entire procedure.  Integer frequencies and finite offending
collections can be chosen in a fixed enumeration.  Polarization
uses measurable bounded signs of the moment functions, and all
membership probabilities are computed from their channel laws.
Thus all resulting input sets are measurable.  No joint choice over
different dyadic grids is required.

\subsubsection*{The spatial likelihood and membership costs}
At a positive-mass discrete state with density $a$ on $p$,
\begin{equation}
 \mathbf P(\text{state})=
       \frac{\mathbf U(p)m}{m_{\rm top}},\qquad
 \mathcal L(Y\mid\text{state})=
       \frac{a}{m}\,\mathcal L(\mathbf U\mid p).
 \label{eq:06-state-law}
\end{equation}
This follows directly by restricting the initial density to the
recorded spatial interval and membership sets.  Define, at each
such state,
\[
 L_{\rm sp}=\log(m/m_{\rm top})+B_{\rm mem}.
\]
A membership outcome with probability $q>0$ replaces $m$ by $qm$
and increases $B_{\rm mem}$ by $-\log q$, so $L_{\rm sp}$ is
unchanged.  At a spatial step its increment is
$\log(2q)$, where $q$ is the conditional probability of the actual
child under the unlabelled true law.  It follows that
\[
 \mathbb E_{\mathbf P}\!\left[
       e^{-L_{\rm sp,new}}\mid\text{past}\right]
       \le e^{-L_{\rm sp,old}}.
\]
Indeed the conditional sum over positive-probability children is
$\sum_q q/(2q)\le1$; equality need not hold if a child has zero
mass.  Thus $e^{-L_{\rm sp}}$ is a nonnegative supermartingale
starting at one.  It remains so after stopping or after keeping the
piece fixed and revealing further spatial bits.

A success has surprisal at most
$C(1+\log(1/b))$ because $\delta\ge cb^2$; a cheap success has
zero surprisal.  The deterministic $k_b$ budget and
\[
 1+\log(1/b)\le C(h/d)k_b
\]
bound the total success surprisal by $Ch$ on every path.
Between one success and the next, let $F$ accumulate only failure
surprisal.  Until that stretch ends, $e^F$ is unchanged at spatial
steps; at a membership test keep it multiplied by $(1-\delta)^{-1}$
on failure and send it to zero on success.  The resulting process is
a nonnegative supermartingale starting at one.  Therefore its
maximal inequality implies, conditionally on the stretch's start,
\[
 \mathbf P(F_{\rm final}>s\mid\text{start})\le e^{-s},
       \qquad s\ge0.
\]
Here $F_{\rm final}$ includes the failures preceding a terminal
success, even though the auxiliary process is then sent to zero.
There are at most $Ch+1$ stretches.  Their conditional second
moments are at most $2$; pad with zero stretches and use the triangle
inequality in $L^2$.  This proves $\|B_{\rm mem}\|_2\le Ch$.
The same maximal inequality applied to $e^{-L_{\rm sp}}$ bounds
the $L^2$ norm of $\max(-L_{\rm sp})_+$ by an absolute constant.
Since $m_{\rm top}\ge1/h$,
\[
 X_*\le\log h+B_{\rm mem}+\max(-L_{\rm sp})_+,
\]
which proves \Cref{eq:source-12}, including the mean at transfer.

We shall also need the expected number $N_{\rm att}$ of regular
attempts.  Each has conditional success probability at least
$c(b_0/8)^2\ge c'h^{-32}$, whereas there are at most $Ch$
successes.  Summing conditional success probabilities over the
finite processing tree gives
\begin{equation}
 \mathbb E_{\mathbf P}N_{\rm att}\le C h^{33}.
 \label{eq:06-attempt-count}
\end{equation}

\subsubsection*{Likelihood evidence and the transfer probability}
The small-mass flag implies $X_*>d$, so its probability is at most
$Ch^2/d^2=o(h^{-10})$.  The count flag is already accounted for by
the charge flag.  To control the latter, attach all the main and
evidence observations specified above.  Compare the true experiment
with a reference law $\mathbf Q$ in which the prior for $Y$ is
$\mathbf U$, the spatial bits and all channel draws still observe
$Y$, but each membership outcome is an independent toss with the
state's prescribed probabilities $\delta,1-\delta$.  On a recorded
branch those probabilities are constants in $Y$.  All subsequent
states and channel parameters are the same predetermined functions
of its spatial and membership history as in the true experiment.
One may stop any zero-true-mass branch in the reference as well.

Let $\mathcal O$ denote the terminal observations through live
termination, including draws at its final success, and let
$R=d\mathbf P_{\mathcal O}/d\mathbf Q_{\mathcal O}$.  On true
histories,
\begin{equation}
 R\le\exp(d+\log h+B_{\rm mem}).
 \label{eq:06-reference-ratio}
\end{equation}
To verify this even for histories of different lengths, fix one
complete discrete history and all its label values, using densities
for the latter.  The true likelihood integrates the same channel
weights over inputs compatible with both the spatial bits and the
actual membership sets, with prior density at most
$e^d/m_{\rm top}$.  The reference integrates over all inputs
compatible with the spatial bits, and its additional coin factors
have product $e^{-B_{\rm mem}}$.  The domain of the true integral
is a subset of that reference domain.  This proves
\Cref{eq:06-reference-ratio} on each terminal branch separately.

At every regular attempt whose observed past is admissible, multiply
a score by the following factor, writing $E_{\rm ev}$ for the
evidence event in \Cref{eq:source-10}:
\[
 \frac{\exp\bigl((t/16)\mathbf1_{
       \{\mathrm{success}\}\cap E_{\rm ev}}\bigr)}
      {1+e^{-t/16}};
\]
otherwise leave the score unchanged.  Group the fresh observations
of this test into that step.  The conditional reference probability
of the evidence given success is at most $e^{-t/8}$ by
\Cref{prop:thermal-evidence}.  Consequently the conditional mean
of the numerator is at most
$1+(e^{t/16}-1)e^{-t/8}\le1+e^{-t/16}$.
The score $\mathcal S$ is therefore a nonnegative reference
supermartingale, with $\mathbb E_{\mathbf Q}\mathcal S\le1$ at
termination.  Change of measure gives the particularly useful bound
\begin{equation}
 \mathbf P(\mathcal S>Re^d)
   =\int_{\{\mathcal S>Re^d\}}R\,d\mathbf Q
   \le e^{-d}\mathbb E_{\mathbf Q}\mathcal S\le e^{-d}.
 \label{eq:06-score-transfer}
\end{equation}

The simultaneous true admissibility and evidence assertion of
\Cref{prop:thermal-evidence} fails with probability $o(h^{-20})$
under the count cap; it includes evidence at a success that
immediately transfers.  Also
$\mathbf P(B_{\rm mem}>d)\le Ch^2/d^2$ by
\Cref{eq:source-12}.  Finally, the total logarithmic denominator
cost in the score is at most
$N_{\rm att}e^{-c d/h}$.  Its expectation tends to zero faster than
$h^{-10}$ by \Cref{eq:06-attempt-count}, so with that allowed
exception it is at most one.  Outside these exceptions and
\Cref{eq:06-score-transfer},
\[
 \frac{\mathcal Q}{16}-1
       \le\log\mathcal S
       \le\log R+d
       \le2d+\log h+B_{\rm mem}
       \le3d+\log h.
\]
This excludes $\mathcal Q>1000d$ for large $d$.  It proves the
stated transfer probability.  The score and evidence draws have
only served to estimate this probability; they never select a
piece or an update.

\subsubsection*{The sum of live plateau weights}
At a plateau's root its mean is $\alpha_e$, so
\Cref{eq:06-state-law} gives
\[
 \sum_{e\text{ live}}\mathbf U(p_e)\alpha_e
       =m_{\rm top}\,
          \mathbb E_{\mathbf P}[\text{number of live plateau starts}].
\]
Epoch starts number at most $Ch+1$ on each path.  It remains to
bound upward resets without multiplying that bound by another $h$.

Consider one epoch, with start and end stopping times $\sigma,\tau$
in the unlabelled filtration.  Include its ending membership success
in $\tau$ if there is one.  Write
$\Delta=L_{\rm sp}(\tau)-L_{\rm sp}(\sigma)$, and let $N$ be its
number of upward resets.  At its $j$th reset time $\rho_j$,
\[
 L_{\rm sp}(\rho_j)\ge L_{\rm sp}(\sigma)+j\log2:
\]
the normalizations have more than doubled at each reset and the
membership surprisal is nondecreasing.  For any fixed $r>0$,
optional sampling of the inverse likelihood gives
\[
 \mathbf P\bigl(L_{\rm sp}(\tau)<L_{\rm sp}(\rho_j)-r
                 \mid\mathcal F_{\rho_j}^{\rm disc}\bigr)\le e^{-r},
 \qquad
 \mathbb E[(-\Delta)_+\mid\mathcal F_\sigma^{\rm disc}]\le1.
\]
Here $\mathcal F^{\rm disc}$ denotes the spatial and membership
history, not the uniform-output filtration used earlier.  The
second bound follows by integrating the corresponding exponential
tail.  Every reset not in the exceptional event in the first bound
has $j\log2\le\Delta+r$.  Summing reset indicators and then
conditioning at the epoch start yields
\[
 (1-e^{-r})\mathbb E[N\mid\mathcal F_\sigma^{\rm disc}]
 \le\frac{\mathbb E[(\Delta+r)_+
                       \mid\mathcal F_\sigma^{\rm disc}]}{\log2}
 \le\frac{\mathbb E[\Delta\mid\mathcal F_\sigma^{\rm disc}]+r+1}
              {\log2}.
\]
Sum this inequality over epochs, weighting by the probability that
each epoch starts, and set $r=2$.  From one epoch's end to the next
epoch's start $L_{\rm sp}$ does not change, because only a membership
step separates them.  Thus all $\Delta$ terms telescope to the
terminal $L_{\rm sp}$.  At termination,
\[
 m\le2e^{4h},\qquad
 L_{\rm sp}\le4h+\log2+\log h+B_{\rm mem}.
\]
The factor two allows the last spatial step that first crosses the
common cutoff; membership restrictions cannot increase the mean.
The expected total number of resets is therefore $O(h)$ by
\Cref{eq:source-12}.  Since $m_{\rm top}\le2$, this proves the
plateau weight bound.

For a fixed plateau, the portions removed from its survivor before
its end are disjoint subsets of its original input mass.  Multiplying
each local removed mean by $\mathbf U_e(p)$ and summing counts each
removed input at most once.  Division by the entry mean $\alpha_e$
gives $\mathbb E_{\mathbf U_e}\sum\xi_p\le1$.  The identical argument
from a final post state $p_0$, excluding its past removals, starts
with mass $m_{p_0}/\alpha_e$.  This proves the remaining assertions
in part \textup{(iii)}.

\subsubsection*{The information cost of main labels}
We first recall the finite entropy accounting used here.  For any
finite remaining spatial bit string, track its conditional relative
entropy to uniform future bits.  Revealing the next bit consumes,
in conditional expectation, precisely that bit's relative entropy
to fair, by the chain rule.  An intervening observation increases
the expected deficit by its conditional information about the
remaining string, at most its conditional information about $Y$.
The remaining deficit is nonnegative, so the same inequality holds
when the process is stopped.  The initial deficit is at most
\[
 D(\mathbf P\|\mathbf U)
   =\mathbb E_{\mathbf P}\log(\rho/m_{\rm top})
   \le d+\log h.
\]
For a membership bit its information cost is at most its conditional
entropy.  Averaging and removing all observed-label conditioning
only increases this entropy, so the total membership cost is at
most $\mathbb E B_{\rm mem}=O(h)$.

For the main-label filtration we retain all main observations, but
not the evidence clones.  The initial type-$0$ observation has zero
information.  At a success at $p$, let $Z$ be the old main
observation and $V$ the new observation of type $A^+$.  Given the
full main-observation history through that success, compare the
conditional channel of $V$ with the probability density
\[
 Q_Z(v)=\frac{\exp[-A(v-Z)-S_p(v-c_p)]}
              {\int\exp[-A(w-Z)-S_p(w-c_p)]\,dw},
       \qquad c_p=\operatorname{center}(p).
\]
For every conditional prior on $Y$, its information about $V$ is
bounded above by the expected relative entropy of the channel to
any fixed comparison density.  The centered envelope bounds the
logarithm of the denominator above by $-F(A+S_p)$.  Expanding that
relative entropy and dropping the nonpositive new-form energy gives
\begin{equation}
 \begin{aligned}
 \mathrm I(Y;V\mid\text{history through success})
 \le{}& F(A^+)-F(A+S_p)\\
 &+\mathbb E[A(V-Z)+S_p(V-c_p)
                         \mid\text{history through success}].
 \end{aligned}
 \label{eq:06-channel-information}
\end{equation}
The notation denotes the usual averaged conditional-information
bound; it may first be read for each fixed value of the history.

Now average this inequality over the main-observation histories at
the fixed discrete success state.  Conditional on $Y$ and that
discrete state, the old and new errors still have their prescribed
independent laws, because every membership decision is input-only.
The phase triangle inequality and form monotonicity imply
\[
 \mathbb E A(V-Z)
   \le2\int A K_{A^+}+2\int A K_A
   \le4\int A K_A
   \le C h(dh^{-400}+1).
\]
This averaging step is essential: conditioning on an observed old
label alone would not give its unconditional error law.  Also
$Y\in p$ and $A^+$ contains $S_p$, so
\[
 \mathbb E S_p(V-c_p)
    \le2\mathbb E S_p(V-Y)+2S_p(Y-c_p)\le C.
\]
Multiply each such estimate by the true probability of its success
state and sum.  Those probabilities sum to the expected number of
successes, at most $Ch$.  Their total old-error energy is therefore
\[
 O\bigl(h^2(dh^{-400}+1)\bigr)=o(d),
\]
and their spatial energy is $O(h)$.  The relative metric terms in
\Cref{eq:06-channel-information} sum to $O(d)$ on every path by
part \textup{(i)}.  Thus all incremental main-channel informations
together cost $O(d)$, including a terminal cap-crossing success.

\subsubsection*{The Fisher sum and the unlabelled deficits}
At a final live state $p$, the conditional density of the current
main label, ignoring the older labels, is $\alpha_e M/m_p$.  Given
this label, the next spatial bit has bias $D/M$.  Hence its expected
squared bias at this state is
\[
 \frac{\alpha_e}{m_p}\int\frac{|D|^2}{M}
       =\frac{\alpha_e}{m_p}J_p.
\]
Multiplication by the state's probability in
\Cref{eq:06-state-law} gives exactly
$\mathbf U(p)\alpha_e J_p/m_{\rm top}$.  Conditioning also on
all older main labels can only increase the average squared
conditional bias, by conditional Jensen.  For a binary law of bias
$b$, its relative entropy to a fair bit is at least $b^2/2$.
The finite entropy accounting just proved therefore bounds the sum
of these weighted $J_p$'s by a constant times
$m_{\rm top}(d+\log h+\mathbb E B_{\rm mem}+O(d))\le Cd$.
This proves \Cref{eq:source-13} with no lower bound on the surviving
mass needed in this conversion.

Finally ignore all labels.  The conditional expected increment of
$L_{\rm sp}$ at a spatial bit is exactly its relative entropy to
fair, while memberships leave $L_{\rm sp}$ unchanged.  Starting
from $L_{\rm sp}=0$, the expected sum of bit deficits equals its
expected terminal value.  On any finite spatial extension with the
last piece fixed, $m\le e^d$, so
\[
 L_{\rm sp,final}\le d+\log h+B_{\rm mem}.
\]
If that extension stops at the common density cutoff, instead
$m\le2e^{4h}$ and
$L_{\rm sp,final}\le4h+\log2+\log h+B_{\rm mem}$.
Taking expectations and using \Cref{eq:source-12} gives respectively
$O(d)$ and $O(h)$.  This completes the proof.
\end{proof}
